\section{Markings and a constant finite normalization}
\label{sec:marking}

The geometric constancy in Proposition~\ref{ro:constancy} concerns the
fibers of $x$ along a geometric fiber of $h$.  We next control the base of
that family.  The purpose of the markings below is to recover a
finite cover of the geometric generic fiber of $Z\to T_0$, together
with the divisorial images of the old boundary.  The markings will then be discarded. There are two separate
conclusions: the finite normalization and its morphism must descend,
and each positive old-boundary image must become fixed. The second
conclusion supplies the actual multiplicity-one component needed in
the ramification argument of Section~\ref{ds:section}.

\subsection{The geometric generic parameter space}

Fix the algebraic closure $\Omega$ of $\C(Y)$ used to define
$\Var(f)$, and put
\[
 b=\overline{\C(T_0)}\subset\Omega,
\]
where the bar denotes the algebraic closure inside $\Omega$.

\begin{lemma}\label{mk:setup}
Let $H_0,I,P_0$ be the geometric generic fibers over $T_0$ of $W,S,Z$,
respectively.  Write $B_0$ and $D_I$ for the restrictions of $B$ and
$D_S$.  These are smooth projective varieties over $b$, with the indicated
SNC boundary supports.  The induced morphisms
\[
 p_I:H_0\longrightarrow I,\qquad a_0:H_0\longrightarrow P_0
\]
have geometrically integral generic fibers, and the combined morphism
satisfies
\begin{equation}\label{mk:generic-diagram}
 H_0\longrightarrow I\times_b P_0
 \quad\text{dominant and generically finite},\qquad
 \dim P_0=\dim(H_0/I)=d.
\end{equation}
Moreover,
\begin{equation}\label{mk:log-base}
 \kappa(I,K_I+D_I)\geq0,
 \qquad B_0\geq(p_I^*D_I)_{\mathrm{red}},
\end{equation}
and a geometric general fiber $H_{0,t}$ of $a_0$ satisfies
\begin{equation}\label{mk:zero-fiber}
 \kappa(H_{0,t},K_{H_{0,t}}+B_0|_{H_{0,t}})=0.
\end{equation}
\end{lemma}

\begin{proof}
The diagram and dimension assertions follow from
Propositions~\ref{it:diagram} and~\ref{it:dimensions}. The latter
identifies $\C(S)=\C(Y)$ inside $\C(X)$, so the present parameter
field is the embedded field $b(I)=b\,\C(Y)$ of
Remark~\ref{it:constant-field-interface}. The generic extensions for
$W\to S$ and $W\to Z$ are regular. Restriction to the generic point
of $T_0$ and extension to $b$ therefore preserve geometric integrality. Stein
factorization over the normal bases also gives connected fibers.  Generic
smoothness, applied also to the finitely many boundary strata, gives the
stated smoothness and SNC properties.  The first assertion of
\eqref{mk:log-base} is the geometric generic form of the corresponding
assertion for $S_t$ in Proposition~\ref{it:dimensions}; nonzero divisible
sections are preserved by field extension.  The second follows from
Proposition~\ref{ro:divisors}.  Finally, \eqref{mk:zero-fiber} is precisely
Proposition~\ref{ro:trivial} after this generic restriction.
\end{proof}

Proposition~\ref{num:model-pair} applies over the algebraically closed
field $b$ of characteristic zero. Its base hypothesis is precisely
$\kappa(I,K_I+D_I)\geq0$, as supplied by Lemma~\ref{mk:setup}.
In particular, the base pair in the marked-model argument is $(I,D_I)$.

\subsection{Recognizable markings}

The generic $I$-fiber of $H_0$ is generically finite over $P_0$.
We will first decrease the positive horizontal boundary coefficients
to make its pair klt, and then add pullbacks of divisors from $P_0$
to make its adjoint divisor big. The marking line bundles restrict
trivially to general $a_0$-fibers, while decreasing the old boundary
can only lower its adjoint dimension from zero. Upper addition
therefore bounds the total section dimension by $d$, and
Proposition~\ref{num:model-pair} bounds it below by $d$ plus the
variation of the marked model pair. The marked model must become
constant after a finite parameter extension. We choose its markings
so that they also recover the map to $P_0$, and hence the normalization
of $P_0$ in the model's function field. Once this normalization is fixed,
Proposition~\ref{ro:constancy} will supply a reference variety for
the $J$-family over its function field.

\begin{proposition}\label{mk:product}
After a finite extension of $b(I)$ there is a normal projective variety
$V'$ over $b$, finite over $P_0$, such that the generic $I$-fiber of
$H_0$ is birational over $P_0$ to the base extension of $V'$.  Equivalently,
after replacing $I$ by a smooth projective model of that finite extension,
there is a birational product description over $I\times_bP_0$,
\begin{equation}\label{mk:product-field}
 H_0\dashrightarrow I\times_b V'.
\end{equation}
Here $H_0$ denotes the main component after the parameter extension.
Every component of the old $B_0$ which has a divisorial image on $V'$
over the generic point of $I$ has an image defined over $b$.

Put
\begin{equation}\label{mk:fields}
 k=b(V'),\qquad K=k(I).
\end{equation}
The original generic $x$-fiber, viewed by \eqref{mk:product-field} as a
variety $J/K$, is geometrically birational to a smooth projective
geometrically integral variety $J_0/k$.  One has
$\kappa(J_{0,\bar k})=0$, and, for the minimal positive pluricanonical
index $p$ of $J$, there is a nonzero ordinary $p$-pluricanonical form on
$J_0$ and
\begin{equation}\label{mk:reference-form}
 \dim_k H^0(J_0,pK_{J_0})
 =\dim_{\bar k}H^0(J_{0,\bar k},pK_{J_{0,\bar k}})=1.
\end{equation}
\end{proposition}

\begin{proof}
If $d=0$, the generic fiber of $p_I$ is geometrically integral and
zero-dimensional, so its degree is one.  The variety $P_0$ is a point
and the product assertion holds with $V'=\operatorname{Spec}b$.
The construction of $J_0$ at the end of the proof applies in this case
as well.  Assume for the marking construction that $d>0$.

\smallskip\noindent\textbf{The relative model.}\enspace
Let $E_0=b(I)$ and $H_\eta=(H_0)_{E_0}$.  Decrease every positive
coefficient of $B_0$ horizontal over $I$ slightly, keeping it positive
and strictly less than one; leave the vertical coefficients unchanged.
Denote the resulting boundary by $B_0^-$.  Since the generic support is
SNC, the pair $(H_\eta,B_\eta^-)$ is klt.  The morphism
\[
 a_\eta:H_\eta\longrightarrow (P_0)_{E_0}
\]
is generically finite.  Consequently $K_{H_\eta}+B_\eta^-$ is big
relative to $a_\eta$, since its generic fiber is zero-dimensional.
There is therefore a relative log canonical model
\[
 H_\eta\dashrightarrow C_1\xrightarrow{\,c\,}(P_0)_{E_0}
\]
by \cite[Theorem~1.2(2)]{BCHM}.  To apply that theorem over an
algebraically closed ground field, one can first shrink $I$ so that the
spread of the decreased pair is klt, apply the theorem to the projective
morphism to $I\times_bP_0$, and then restrict to its generic $I$-fiber.
The characteristic-zero field comparison in
Section~\ref{num:comparisons} justifies this use of the cited complex
formulation over $b$. The model $C_1$ is normal, and its adjoint divisor is relatively ample.
The canonical rational map is a birational contraction, so it extracts
no divisors; see \cite[Definition~3.6.6]{BCHM}.  The map $c$ is proper,
surjective and generically finite, and factors through the normalization
of $(P_0)_{E_0}$ in $E_0(H_\eta)$.

\smallskip\noindent\textbf{Choice of the markings.}\enspace
Choose a very ample line bundle $L$ on $P_0$ and a general basis
$x_0,\ldots,x_N$ of $H^0(P_0,L)$. The coordinate divisors will
recover each ratio $x_j/x_0$ up to a scalar. The divisor of
$x_0+x_j$ will determine that scalar, so that the marked model
remembers the morphism to this fixed $P_0$. We therefore use
the divisors
\begin{equation}\label{mk:markings}
 (x_0=0),\qquad (x_j=0),\qquad (x_0+x_j=0)
 \quad (1\leq j\leq N).
\end{equation}
Their pullbacks receive distinct positive rational coefficients, chosen
sufficiently small that the resulting generic pair is still klt.  They
are also chosen distinct from every coefficient of the decreased generic
strict boundary $B_\eta^-$.  The finite number of dependencies in
\eqref{mk:markings} is harmless: on a fixed log resolution, sufficiently
small coefficients preserve all the strict klt discrepancy inequalities.
Add sufficiently many further general members of $|L|$, again with
small coefficients less than one and with no coefficient equal to one
of the displayed labels.  These can have arbitrarily large total degree.
General members of the free pullback system meet the strata on a
resolution transversely, so the generic pair remains klt.

Make the following avoidance choices at the same time.  None of the
hyperplanes contains the image of any relevant divisorial exceptional
component on the fixed generic models, any old generic boundary
component, or any branch divisor.  In addition, none contains the
$P_0$-image of an $I$-vertical divisor of $H_0$ having proper image in
$P_0$.  There are only finitely many divisors in the last condition:
such a divisor has image of codimension at least two under the
generically finite morphism $H_0\to I\times_bP_0$, so it is one of its
divisorial exceptional components.  An $I$-vertical divisor dominating
$P_0$ cannot be a component of the pullback of a hyperplane.  Thus the
last condition prevents any added marking from overlapping a vertical
old coefficient-one component or acquiring a vertical divisorial
multiplicity.  It also makes the corresponding total boundary a
boundary with coefficients at most one.

These finitely many Bertini, avoidance, and klt conditions define a
nonempty open in the parameter space of the chosen sections, initially
over $E_0$. Lemma~\ref{prelim:generic}, applied after extending $E_0$
to its algebraic closure, shows that this open contains a $b$-point.
Thus all the markings can be chosen over $b$. The same argument applies
to the open set of bases of $H^0(P_0,L)$ and to each additional member.

Write $A_{\mathrm{mark}}$ for the weighted sum of all the chosen
divisors on $P_0$, and set
\[
 \Delta=B_0^-+a_0^*A_{\mathrm{mark}}
 \quad\text{on }H_0.
\]
The total degree of $A_{\mathrm{mark}}$ is chosen sufficiently large
that
\[
 K_{C_1}+(B_\eta^-)_{{C_1}}+c^*A_{\mathrm{mark}}
\]
is ample: a relatively ample divisor plus a sufficiently large
pullback of an ample divisor is ample.  Twisting by pullback from
$P_0$ tensors each divisible graded piece of the relative adjoint ring
by the corresponding power of a line bundle on $P_0$.  It therefore
leaves the relative $\operatorname{Proj}$, and hence the model, unchanged.
Equivalently, on a common resolution the exceptional comparison for
the old adjoint divisor remains identical after adding this pullback
on both sides.  Hence $C_1$ is now the absolute log canonical model of
$(H_\eta,\Delta_\eta)$.

The pullbacks of the displayed divisors on $C_1$ are their complete
pullback divisors, each with multiplicity one at every component.  A
divisorial exceptional component occurs in a pullback only when its
image is contained in the hyperplane; merely meeting its image does not
produce that component.  The avoidance of the branch divisors gives
multiplicity one at the other generic points.  No old boundary prime
is shared with a marking.  Thus their distinct coefficients identify
the complete displayed divisors on the canonical model, even when a
pullback is reducible.

\smallskip\noindent\textbf{The upper bound and the model-pair estimate.}\enspace
On a general $a_0$-fiber the marking line bundles restrict trivially,
and $B_0^-\leq B_0$.  Lemma~\ref{mk:setup} therefore gives restricted
adjoint dimension at most zero.  The upper addition estimate
in Lemma~\ref{prelim:addition} gives
\begin{equation}\label{mk:upper}
 \kappa(H_0,K_{H_0}+\Delta)\leq\dim P_0=d.
\end{equation}
If the restricted section system is empty, the same bound follows
because the total system is empty.  Only the upper bound is needed.

Resolve the support of $\Delta$, writing
$\mu:\widehat H\to H_0$.  Over the generic point of $I$, use the
crepant boundary and replace its negative exceptional coefficients by
zero; keep all strict coefficients.  Call the resulting generic
boundary $\widehat\Delta_\eta$.  The generic pair is effective and klt,
and
\begin{equation}\label{mk:effective-crepant}
 K_{\widehat H_\eta}+\widehat\Delta_\eta
 =\mu_\eta^*(K_{H_\eta}+\Delta_\eta)+E_{\mathrm{exc}},
 \qquad E_{\mathrm{exc}}\geq0
 \quad\text{$\mu_\eta$-exceptional}.
\end{equation}
For divisible $m$, normality and testing at codimension-one points give
$\mu_{\eta*}\cO(mE_{\mathrm{exc}})=\cO$.  The generic adjoint rings are
therefore identical.  The generic model pair remains the marked pair
on $C_1$: exceptional divisors added in \eqref{mk:effective-crepant}
are contracted to the original generic source, and the canonical map
from that source extracts no divisors.

On $\widehat H$, assign coefficient one to each new exceptional divisor
vertical over $I$, and retain the specified coefficients on all other
components.  This gives an effective rational SNC boundary
$\widehat\Delta$, with generic restriction just described.  All strict
components above $D_I$ already had coefficient one, and any newly
created component above it is vertical over $I$.  Consequently
\begin{equation}\label{mk:inverse-boundary}
 \widehat\Delta\geq
 \bigl((p_I\mu)^*D_I\bigr)_{\mathrm{red}}.
\end{equation}
The total adjoint sections on $\widehat H$ inject into those on $H_0$:
a rational pluriform that is a section upstairs passes the unchanged
strict-divisor tests at every prime downstairs.  This argument does not
require the vertical exceptional coefficients to be crepant.  In
particular, even if one of these exceptional divisors dominates $P_0$,
\eqref{mk:upper} still yields
\begin{equation}\label{mk:resolved-upper}
 \kappa(\widehat H,K_{\widehat H}+\widehat\Delta)\leq d.
\end{equation}

The markings now define a big generic model pair, while the total
section dimension still has the upper bound $d$. We apply
Proposition~\ref{num:model-pair} to
\[
 p_I\mu:(\widehat H,\widehat\Delta)\longrightarrow(I,D_I).
\]
Both varieties are smooth projective over $b$; the morphism is
surjective with connected fibers and relative dimension $d$; the
boundary is effective rational SNC; \eqref{mk:inverse-boundary} holds;
and the base has nonnegative log Kodaira dimension by
\eqref{mk:log-base}.  Its geometric generic pair is klt and has big
adjoint divisor by the construction of the absolute ample model $C_1$.
Thus Proposition~\ref{num:model-pair} and \eqref{mk:resolved-upper} give
\[
 d+\Var(p_I\mu,\widehat\Delta_{\mathrm{hor}})
 \leq\kappa(\widehat H,K_{\widehat H}+\widehat\Delta)\leq d.
\]
The variation here is exactly that of the generic log canonical model
\emph{pair}. It is zero. The finite-extension conclusion of
Proposition~\ref{num:model-pair} makes that model pair isomorphic, after a finite
extension $E/E_0$, to the base extension of a model pair over $b$.
Denote its underlying normal projective variety by $C_*$.  Its
coefficient-labeled displayed divisors are now defined over $b$.

\smallskip\noindent\textbf{Recovering the map to $P_0$.}\enspace
View $C_1$ after extension to $E$ as $(C_*)_E$.  The displayed markings
fix both
\[
 \operatorname{div}(x_j/x_0)
 \quad\text{and}\quad
 \operatorname{div}(1+x_j/x_0).
\]
These are differences of the corresponding complete marking divisors.
In particular, the first is the base extension of a divisor over $b$
which is principal after extension to $E$.  Cartierness descends under
this faithfully flat extension.  Triviality of its associated line
bundle descends as well: its space of sections commutes with field
extension, and a nonzero descended section becomes a scalar multiple
of a nowhere-vanishing section upstairs.  We may therefore write
\begin{equation}\label{mk:coordinate-functions}
 x_j/x_0=c_jf_j,
 \qquad f_j\in b(C_*),\qquad c_j\in E^*.
\end{equation}
Here the quotient of two functions with the same divisor is a constant,
since $C_*$ is proper and geometrically integral.

The sum marking $(x_0+x_j=0)$ has a $b$-point off the two coordinate
markings.  Indeed, for a general basis this is a nonempty open part of
the pullback of that hyperplane, by surjectivity onto $P_0$ and $d>0$;
as it is defined over the algebraically closed field $b$, it has a
$b$-point.  Evaluation at such a point in
$1+c_jf_j=0$ gives $c_j\in b$.  Thus every coordinate ratio of the
map to the fixed projective embedding of $P_0$ belongs to $b(C_*)$.
The map itself descends to a morphism $C_*\to P_0$, either by descent
of its graph or by descent of regularity of the resulting rational
map.

Let $V'$ be the normalization of $P_0$ in the finite extension
$b(C_*)/b(P_0)$.  It is finite over $P_0$.  The map $C_*\to P_0$
factors properly and birationally through $V'$.  Normalization commutes
with the present field extensions: $b$ has characteristic zero and is
algebraically closed, so $V'$ is geometrically normal; its base extension
is finite normal with the required fraction field and hence has the
normalization's universal property.  The original generic source
$H_\eta$ after extension to $E$ also maps properly and birationally to
$V'_E$.  This proves \eqref{mk:product-field}.

A proper birational morphism to a normal variety is an isomorphism at
the generic point of every target divisor.  Consequently a prime on
$V'_E$ has its unique strict prime both on the original generic source
and on $(C_*)_E$.  If the old boundary coefficient at that prime is
positive, its decreased coefficient is still positive and occurs in
the constant model pair.  Its image on $V'$ is therefore defined over
$b$.  This proves the asserted control of old boundary images.  Notice
that this reasoning uses the old boundary; a newly assigned
exceptional coefficient on $\widehat H$ cannot hide such a divisor.

\smallskip\noindent\textbf{Choosing the reference fiber.}\enspace
Replace $I$ by a smooth projective model of $E$, and use the notation
\eqref{mk:fields}.  All subsequent parameter extensions will be made
relative to this choice.  The original $x$-fiber gives a smooth
projective geometrically integral variety $J/K$ with geometric Kodaira
dimension zero.  Proposition~\ref{ro:constancy} applies after extending
$k$ to an algebraic closure: the $I$-directions in
\eqref{mk:product-field} lie over one geometric generic $Z$-point.
Thus this family is geometrically birationally constant in those
directions.

Spread $J$ as a smooth projective family over a nonempty open of $I_k$.
Fix one geometric birational identification with a reference variety
over $\bar k$, together with its inverse. Their finitely many
coefficients lie in a finite extension of $\bar k(I)$. Normalize an
affine open of $I_{\bar k}$ in that extension. After shrinking, this
is a finite surjective cover on which the maps and their inverse
identities spread. Shrink again so that their fiberwise domains are
dense and the two maps are inverse there. This uses one chosen
birational identification and finitely many equations, not a
finite-type parameter space for all birational maps.

The image of this comparison cover contains a nonempty open of
$I_{\bar k}$. Intersect it with the smooth family open. By
Lemma~\ref{prelim:generic}, that intersection contains a point of
$I(b)$; a point of the finite comparison cover above it exists over
$\bar k$. Specializing the original $k$-family at this $b$-point
gives a smooth projective variety $J_0/k$. Its base extension to
$\bar k$ is birational to the chosen reference, hence geometrically
birational to $J$. Geometric
integrality and geometric Kodaira dimension zero follow.  Ordinary
birational invariance of pluricanonical sections and their compatibility
with field extension give \eqref{mk:reference-form}; in particular a
nonzero ordinary $p$-pluriform is defined over $k$ itself.
\end{proof}

\subsection{A multiplicity-one divisor at every moving test}

The boundary control has a stronger consequence than the existence of
some reduced component in a degeneration.  It supplies a reduced
component attaining the precise threshold used in the relative-form
comparison.

\begin{definition}\label{mk:moving-definition}
In the notation \eqref{mk:fields}, let $E/b(I)$ be a finite extension.
A \emph{moving test} is a prime divisor $P$ of $V'_E$ whose image in
$V'$ is dense.  Its valuation is trivial both on $k=b(V')$ and on $E$.
\end{definition}

For example, if $V'=\mathbf P^1_b$ with coordinate $x$ and
$b(I)=b(t)$, the divisor $x=t$ is moving, whereas $x=c$ for
$c\in b$ is fixed. The next lemma applies to moving divisors over
every finite extension of the parameter field.

\begin{lemma}\label{mk:moving}
At every moving test, including after any finite parameter extension
as in Definition~\ref{mk:moving-definition}, the old coefficient $B_P$
from Proposition~\ref{ro:divisors} is zero.  For the relative ordinary
$p$-pluricanonical generator of the $J$-family, there is an actual source
prime $Q$ above $P$ such that
\begin{equation}\label{mk:threshold-minimizer}
 m_Q=1,
 \qquad
 \frac{1+r_Q}{m_Q}
 =\inf_{Q'\mapsto P}\frac{1+r_{Q'}}{m_{Q'}}.
\end{equation}
Here $m_Q=\ord_Q(x^*P)$ and
$r_Q=p^{-1}\ord_Q(\omega_J)$ are measured in the actual relative
canonical bundle, and the infimum includes the divisorial valuations
on higher smooth source models.  The prime $Q$ is inherited from the
actual source in the rank-one construction after the indicated
parameter extension, not inserted by a new reduced-exceptional
boundary convention.
\end{lemma}

\begin{proof}
A moving test is trivial on $k$ because its residue field contains the
function field of its dense image in $V'$; it is trivial on $E$ because
it is a divisor on an $E$-variety.  The proper birational maps to the
finite normalization used in Proposition~\ref{mk:product} identify its
generic point with a strict divisor of the old generic source $H_0$.
If that divisor had positive old $B_0$ coefficient, the boundary-image
assertion of Proposition~\ref{mk:product} would place its image in a
fixed $b$-divisor on $V'$.  Such a divisor, after any parameter extension,
cannot dominate $V'$.  Thus its old coefficient is zero.

We explain why these are indeed the old tests and orders.  A finite
separable extension of the parameter field is finite \emph{\'etale} as
a field morphism.  Base change is therefore finite \emph{\'etale} on
both the base and the source of the generic family.  Height-one primes
lie over height-one primes, and actual source components persist,
possibly splitting.  Their pullback multiplicities and relative
canonical orders are unchanged.  Generic restriction and adjunction
cancel the parameter-canonical factors in the relative canonical
bundles.  The same assertion for the all-valuation infimum follows by
base changing a fixed log resolution computing the threshold.
Extensions involved in reaching $b=\overline{\C(T_0)}$ are algebraic,
and each prime and the finite data just used descend to a finite stage.
Consequently the unchanged-divisor and \emph{\'etale}-extension
compatibility in Proposition~\ref{ro:divisors} applies at this strict old
prime.  No new exceptional boundary coefficient enters the calculation.

Choose an actual component $Q$ attaining the first minimum $t_P$ of
Proposition~\ref{ro:divisors}.  With $d_Q$ its old source boundary
coefficient, that proposition gives
\[
 B_P=1-\lambda_P+t_P=0,
 \qquad
 1=\lambda_P-t_P
 \leq\frac{1-d_Q}{m_Q}\leq1.
\]
Since $m_Q$ is a positive integer and $0\leq d_Q\leq1$, all these
inequalities are equalities.  Hence $m_Q=1$ and $d_Q=0$.  Equality in
the first inequality says that this same actual component attains
$\lambda_P=\inf_{Q'\mapsto P}(1+r_{Q'})/m_{Q'}$.  This proves
\eqref{mk:threshold-minimizer}.
\end{proof}
